\section{Proposed method: side-specific interval updates}\label{sec:separable}
The first construction is a scalar interval proximal operation inserted
into an original-row quadratic-programming split. Its input is not a
single generic softness number. Each finite lower and upper side has its
own nonnegative linear price and quadratic curvature. This permits a
preferred lower departure to have a different economic or engineering
meaning from an upper departure. Physical hard rows remain a separate
indicator constraint. The construction is related to generalized slack
projection, but its lower and upper branches are derived directly rather
than justified by assuming a sign-symmetric penalty~\cite{gsp,lew}.

For one nonempty interval $[\ell,u]$, write
\begin{equation}\label{eq:paper-side-cost}
 g(z)=a_L(\ell-z)_++a_U(z-u)_+
       +\tfrac12b_L(\ell-z)_+^2+\tfrac12b_U(z-u)_+^2,
 \quad a_L,a_U,b_L,b_U\geq0.
\end{equation}
If a side is absent its term is omitted. A zero linear or quadratic
coefficient is admissible; a negative quadratic coefficient is not. On a
hard row $g$ is replaced by the interval indicator rather than approximated
by a very large finite price. The distinction avoids confusing an elastic
optimum with physical hard feasibility. It also keeps the mathematical
domain unchanged when a local cost multiplier is varied.

For a temporary scalar $v$ and fixed positive metric $h$, define
\begin{equation}\label{eq:paper-scalar-prox}
 p_h(v)=\argmin_z\left\{g(z)+\tfrac h2(z-v)^2\right\}.
\end{equation}
The displacement term is strongly convex even if both quadratic side
prices vanish, so the proximal minimizer is unique. The penalty is zero
on the preferred interval and nondecreasing in each exterior departure.
Consequently an optimizer need not be sought beyond the segment joining
$v$ to its nearest interval endpoint. This geometric restriction is a
consequence of convex monotone exterior prices, not a numerical clipping
heuristic that changes an already computed solution.

If $v$ lies inside the interval, $z=v$ simultaneously minimizes the zero
penalty and the squared displacement. If $v>u$, write $d=z-u\geq0$.
The derivative on the positive exterior branch is
\begin{equation}\label{eq:paper-upper-derivative}
 (h+b_U)d+a_U-h(v-u).
\end{equation}
If this derivative is nonnegative at zero, the endpoint satisfies the
subgradient condition. Otherwise its unique positive root gives the
exterior solution. Reflecting the displacement produces the corresponding
lower-side equation with $a_L,b_L$. The complete map is
\begin{equation}\label{eq:paper-side-prox}
p_h(v)=\begin{cases}
\ell-\dfrac{[h(\ell-v)-a_L]_+}{h+b_L},&v<\ell,\\[6pt]
v,&\ell\leq v\leq u,\\[3pt]
u+\dfrac{[h(v-u)-a_U]_+}{h+b_U},&v>u.
\end{cases}
\end{equation}
All denominators remain positive. A hard row uses interval projection,
and a free soft side returns the unconstrained temporary value on that
side. At an equality preference the lower and upper formulas meet at the
same endpoint with their respective subgradient intervals. No auxiliary
departure variable is needed to evaluate this row operation.

The derivation is a standard convex proximal calculation, consistent with
the framework of Moreau and subsequent proximal methods~\cite{moreau,prox}.
It does not establish that unequal side prices were previously
inexpressible. DAQP's existing side-weight interface and an explicit
finite-side lift already represent this objective~\cite{daqp}. The proposed
computational property is the placement of the admitted update in the
original-row native path and its potential effect on complete-operation
cost. A claim of global priority would require stronger historical
evidence than the construction or its numerical agreement supplies.

An equivalent residual description clarifies sign conventions. Let
$C=[\ell,u]$, $r=z-\Pi_C(z)$ and
\begin{equation}\label{eq:paper-asym-residual}
 \Psi(r)=a_U(r)_++a_L(-r)_+
             +\tfrac12b_U(r)_+^2+\tfrac12b_L(-r)_+^2.
\end{equation}
Then $g(z)=\Psi(z-\Pi_C(z))$. A corrective slack $s$ used in
$z+s\in C$ has the opposite sign and cost $\Psi(-s)$. Forgetting
that reflection exchanges lower and upper prices while preserving an
apparently convex objective. Such a solver can converge successfully to
the wrong policy. The independent evaluator therefore reconstructs
departure from the original row value and evaluates the original
coefficients directly.

Dual consistency supplies another independent check. For finite endpoints
and positive quadratic coefficients, the conjugate is
\begin{equation}\label{eq:paper-side-conjugate}
g^*(y)=\begin{cases}
 uy+[(y-a_U)_+]^2/(2b_U),&y\geq0,\\
 \ell y+[(-y-a_L)_+]^2/(2b_L),&y<0.
\end{cases}
\end{equation}
When the corresponding $b$ is zero, its quadratic conjugate term is
replaced by the appropriate price-cap indicator. An absent endpoint
restricts the corresponding multiplier sign rather than introducing
infinite-times-zero arithmetic. These cases follow by maximizing the
linear dual pairing over each exterior ray~\cite{rockafellar}. The
relationship tests price scaling and sign recovery independently of the
primal proximal formula.

For the full problem, original stationarity and subgradient consistency
require
\begin{equation}\label{eq:paper-side-kkt}
 Px+q+A^\top y=0,\qquad
 Ax=\prox_g^{H_0}(Ax+H_0^{-1}y),
\end{equation}
where $H_0$ is an explicitly selected positive diagonal diagnostic metric
with cost per squared row-unit entries, not the iteration's numerical
penalty. A numeric identity is used only in canonical diagnostic coordinates.
This diagnostic remains well
defined at kinks. A small stationarity residual alone is insufficient:
an unconstrained stationary decision outside a priced interval can have
a wrong row multiplier. Conversely, a nonunique dual need not match an
oracle componentwise when it satisfies the correct original conditions.
The comparison therefore retains objective, hard feasibility and method-
appropriate dual diagnostics as distinct quantities.

The native split retains the original variables and row system. With
positive primal regularization $S$ and positive diagonal row metric $H$,
the graph-proximal calculation involves
\begin{equation}\label{eq:paper-original-system}
 P+S+A^\top H A\succ0.
\end{equation}
Changing side prices changes the proximal map, while changing a numerical
metric can change preparation or factorization. Unchanged dimensions do
not imply that every update is free. The tested prepared workspace admits
only fixed matrices, coefficient structures, hard mask and finite-side
patterns, with the permitted vector updates. It refuses a changed model
structure rather than reusing a factor that belongs to another problem.

Algorithm~\ref{alg:paper-separable} states the executable procedure. It
separates semantic admission, numerical iteration and original recovery.
The separate engine owns the native update and its binding; the product
does not emulate a solver in the browser. The original registered
computation uses its exact recorded source rather than silently replacing
it with a later public wheel. That source distinction is necessary for
reproducible timing even when both installations implement the same
mathematical penalty family.

\begin{algorithm}[htbp]
\caption{Execute side-specific interval updates on the original row system, preserve hard policy, and return original-space quantities for independent objective and feasibility evaluation.}\label{alg:paper-separable}
\begin{algorithmic}[1]
\Require Original convex QP, nonnegative side prices, positive metric
\State Refuse unsupported grouped/Huber prices and invalid row metadata
\State Prepare the original-row native linear system
\While{the numerical termination criterion is not reached}
\State Apply interval projection on hard rows
\State Apply \eqref{eq:paper-side-prox} on admitted soft rows
\State Continue the native primal/dual splitting calculation
\EndWhile
\State Recover $x,y$ and original lower/upper departures
\State Recompute objective, hard error and available dual diagnostics
\end{algorithmic}
\end{algorithm}

Physical covariance is an additional semantic restriction. Under a
positive row conversion $z'=dz$, equivalent prices and duals satisfy
\begin{equation}\label{eq:paper-side-covariance}
 a_L'=a_L/d,\quad a_U'=a_U/d,\quad
 b_L'=b_L/d^2,\quad b_U'=b_U/d^2,\quad y'=y/d.
\end{equation}
The bounds and row matrix must also scale. Changing only the numerical
coefficient or the unit label is not an equivalent problem. The concrete
tests exercise original-space covariance and deliberately wrong units;
they do not establish every floating-point quantity invariant under
arbitrarily extreme scaling. Floating-point conditioning remains a
separate numerical issue~\cite{higham,knight}.

The main implementation advantage sought was reduced setup and solve
work against equivalent lifting. The original evaluation does not
establish that advantage: its complete primary interval crosses zero.
That outcome is consistent with a scalar operation being exact yet
contributing little to total validated-call cost. Formulation construction,
factorization, backend maturity, input checks and recovery can dominate
the saved auxiliaries. The implementation and its derivation remain
useful for explaining side-specific semantics, but the failed prediction
cannot be replaced with a favorable case, a selected dimension or a warm
workspace example after the held-out results are known.

Finally, the convergence interpretation remains bounded. Fixed-parameter
convex splitting under an attained saddle point supports convergence in
its stated domain~\cite{eckstein,admm}. The original solver may use
documented adaptive numerical settings, whose observed successful outputs
are tested directly. The scalar formula does not by itself prove every
adaptive schedule, every numerical infeasibility certificate or a
finite-precision error bound. Keeping that distinction explicit avoids
turning a correct local minimization into a claim about all behavior of
a large native optimization implementation.

The distinction between a finite linear price and exact hard feasibility
can be stated more precisely. Suppose the corresponding hard-constrained
convex problem has an attained primal and dual optimum. A linear price that
dominates an appropriate optimal multiplier can support an exact penalty
representation of its optimum~\cite{kerrigan}. Equality at the threshold
can still allow additional elastic minimizers, so a finite user-selected
coefficient cannot be interpreted as an unconditional physical guarantee.
The product instead treats the hard mask as an independent invariant.
This is why changing a cost factor never changes which original rows
are immutable.

Proximal covariance also requires transforming the numerical metric when
comparing individual row operations. Under the positive row conversion
in\eqref{eq:paper-side-covariance}, the temporary value and bounds scale
with the row, and the equivalent metric scales inversely with its squared
unit factor. Holding that metric fixed instead can change the numerical
trajectory while preserving the original optimization target after full
model conversion. Therefore equality of a single proximal iterate and
equality of the final physical optimum are different claims. The tests
state which object is being compared rather than using a relabelled axis
as evidence of dimensional invariance.

Sparse validation similarly has several layers. A compressed-column
object must have consistent pointers, row indices, shape and vector
lengths. That structural admission does not prove that arbitrary supplied
curvature is positive semidefinite. The mathematical method assumes a
convex model, generated comparison cases have the documented curvature,
and backend/independent numerical failures remain observable. Adding a
small permanent diagonal shift to make another model easier would change
its objective. The original profile therefore records numerical settings
and scientific coefficients separately instead of treating both as
interchangeable solver options.

The implementation comparison also separates timed and independent
lifting. The timed lifted path uses the same pinned native backend
family as the candidate, reducing one obvious implementation mismatch.
The accuracy oracle uses stock OSQP with its own tighter settings.
Substituting that oracle's measured time after seeing a result would
alter the comparator. Likewise, comparing an already prepared candidate
to a reference reconstructed on every repetition would answer a reuse
question different from the registered cold-call hypothesis. All those
boundaries remain explicit in the stored observation fields.

The scalar update's arithmetic is only one part of the full operation.
Original objective evaluation requires recovering row values and pricing
their actual departures. Imported or exported results must also preserve
the base model, side coefficients and hard policy. A smaller sparse
system may coexist with similar validation cost or additional conversion
work. The numerical experiment measures the declared whole call; it
does not time each arithmetic primitive independently. A proposed
explanation that thresholding, validation or sparse ordering dominates
would consequently require additional instrumentation before becoming
a measured causal conclusion.

Finally, the original primary sample has a specific dependence structure.
Seven repetitions on each of twenty-four problems generate the reported
paired observations. Several problems share one evaluation seed, which
is why the confidence procedure resamples seed clusters. Treating every
repeated time as a new independent problem would exaggerate information
and answer a different statistical question. The native implementation's
mathematical admission and the primary experimental denominator are both
preserved, even though they constrain different parts of the resulting
claim.
